\documentclass[a4paper,11pt]{article}
\usepackage{jcappub} % JCAP style
\usepackage[T1]{fontenc}

\title{The Echo Equation in Resonant Medium Cosmology: A Finite-Memory Mechanism for Cosmic Expansion}

\author[a]{Jason Watts}
\affiliation[a]{The Experimentalist Lab}
\emailAdd{theexperimentalistlab@outlook.com}

\abstract{
We present the \emph{Echo Equation}, a finite-memory response framework for cosmic expansion within a resonant-medium cosmology. We derive the governing dynamics, show their reduction to $\Lambda$CDM in a suitable limit, and demonstrate observational viability against key probes (BAO, SN, CMB lensing, structure growth). The model yields distinct, testable signatures---including echo-like modulations in the expansion and growth history---which enable near-term falsification with current and upcoming surveys. We conclude with theoretical consistency checks and predictions.
}

\keywords{cosmology of theories beyond the SM, dark energy theory, modified gravity, cosmological parameters from LSS, gravity in more than four dimensions} % adjust as needed

\arxivnumber{} % optional; add when ready

\begin{document}
\maketitle
\flushbottom

\section{Introduction}
\label{sec:intro}
% Motivation: explanatory mechanism vs. effective parameterization; outline of resonant medium view.
% Brief literature context and prior art; where Echo Equation fits and how it differs.
% Paper roadmap.

\subsection{Motivation and background}
% Summarize tensions/limitations in standard modeling that motivate a finite-memory/resonance approach.

\subsection{This work: contributions}
% Bullet the key contributions: derivation, mapping to observables, fits/predictions, falsifiable signatures.

\section{The Echo Equation: finite-memory cosmological dynamics}
\label{sec:theory}
% Define the resonant medium, memory kernel, state variables.
% Derive the Echo Equation starting from first principles/ansatz.
% Show reduction to \LCDM{} in limiting cases; discuss parameters and physical meaning.

\subsection{Memory kernel and resonance parameters}
% Define parameters (e.g., \tau, \kappa, \omega_0, Q) and their roles.

\subsection{Background evolution}
% Map to H(z), effective w(z), distance measures.

\subsection{Linear growth of structure}
% f\sigma_8 predictions; lensing implications; ISW expectations.

\section{Data, likelihoods, and fitting procedure}
\label{sec:data}
% Describe datasets (BAO, SN, CMB lensing, growth). Covariances and consistency cuts.
% Brief the pipeline: joint likelihood, chi^2, priors, samplers/tools.
% Point to code/data availability (appendix).

\subsection{Datasets}
% BAO (e.g., DESI), SN (e.g., Pantheon+), CMB lensing (Planck), RSD/growth.

\subsection{Methodology}
% Likelihood definition, nuisance handling, fitting details; reproducibility notes.

\section{Results}
\label{sec:results}
% Best-fit parameters and uncertainties; model comparison with \LCDM{}.
% Figures: H(z), w(z), growth, residuals, triangle plots.

\subsection{Background expansion and distances}
% Plots/tables and interpretation.

\subsection{Growth and lensing signatures}
% f\sigma_8, A_L or equivalent; echo-imprint discussion.

\subsection{Model comparison and robustness}
% Information criteria or \Delta\chi^2; sensitivity to priors/cuts.

\section{Predictions and tests}
\label{sec:predictions}
% Concrete, near-term tests; where/when the model deviates most from \LCDM{}.
% Forecasts for DESI/Euclid/LSST; cross-correlations (ISW, g×\kappa).

\section{Theoretical consistency}
\label{sec:consistency}
% Thermodynamics and EFT sanity; early-Universe constraints (BBN, N_eff); backreaction checks.
% Connections to other finite-memory/resonant theories.

\section{Discussion and conclusion}
\label{sec:conclusion}
% Synthesis, implications, limitations; roadmap for future work.

\appendix

\section{Derivation details for the Echo Equation}
\label{app:derivation}
% Expanded steps, proofs, and intermediate lemmas.

\section{Spectral parameters and mapping to cosmological observables}
\label{app:spectral}
% From (\tau,\kappa,\omega_0,...) to (w(z), H(z), f\sigma_8).

\section{Reproducibility: code and data}
\label{app:repro}
% Minimal commands, file manifests, and versioning notes.
% Include SHA256 hashes, software versions, and a pointer to artifacts when public.

\acknowledgments
The author gratefully acknowledges drafting and organizational assistance from the ChatGPT-5 model (OpenAI). No funding or external collaboration supported this work.

\paragraph{Data and code availability} Will be made available upon acceptance or earlier via a public repository; pointers and hashes will be added in App.~\ref{app:repro}.

\paragraph{Conflicts of interest} None declared.

\bibliographystyle{JHEP}
\bibliography{refs}

\end{document}
